\documentclass[12pt]{article}
\usepackage{graphicx} % Required for inserting images
\usepackage{geometry}
\geometry{a4paper, margin=2.5cm}
\usepackage[T1]{fontenc}
\usepackage{palatino}
\usepackage{amsmath, amssymb}
\usepackage{array}
\usepackage{titling}

\setlength{\droptitle}{-5em}
\preauthor{\begin{center}\large\vspace{0.5em}}
\postauthor{\end{center}\vspace{-0.5em}}

\title{
    {\scshape  Part IIB Revision Notes} \\ 
    \vspace{0.5em}
    \textbf{4F7: Statistical Signal and Network Models} \\ 
    Networks}
\author{Qianshuo (Harry) Ye}
\date{April 2026}

\usepackage{parskip}
\begin{document}
\maketitle

\section{Introduction to Networks}
\textbf{Adjacency matrix $A$:} binary valued and symmetric matrix with
$$A_{ij} =
\begin{cases}
1 & \text{if there is an edge from node $i$ to node $j$} \\m
0 & \text{otherwise}
\end{cases}
$$

\textbf{Degree of a node $k_i$:} number of edges incident to the node
$$k_i = \sum_{j=0}^{n-1} A_{ij}$$
The average (mean) degree is $k = \frac{1}{n} \sum_{i=0}^{n-1} k_i = \frac{2m}{n}$.

\textbf{Bipartite graph:} a graph where the nodes can be split into 2 groups, such that all edges fall between groups and non within groups.

\textbf{Components:} a maximal set of nodes that can be reached from on another. An undirected graph is \textit{connected} if it has only one component.

\textbf{Shortest path $s(i,j)$:} the minimum number of edges that must be traversed to reach one node $i$ from the other node $j$. If $i$ and $j$ are directly connected, $s(i,j) = 1$. The shortest path is not defined for nodes in different components. Properties of shortest path length (distance): $s(i,i) = 0$, $s(i, j) = s(j, i)$, and $s(i, j) \leq s(i, k) + s(k, j)$.

\textbf{Diameter:} the longest shortest path between any two nodes in the component.

\section{Network Metrics}
\subsection{Node Importance}
\textbf{Eigenvector centrality:} the importance of a node ($x_i$) is proportional to the sum of the importance of its neighbors.
$$x_i = \alpha \sum_{j=0}^{n-1} A_{ij} x_j$$ For some constant $\alpha$. In matrix form, $\mathbf{x} = \alpha A \mathbf{x}$, which means $\mathbf{x}$ is an eigenvector of $A$. If $x_i \geq 0$, the solution is unique and will correspond to the eigenvector of $A$ with the largest eigenvalue.

\textbf{PageRank:} the importance of a node is the probability that a websurfer is on that node in the long time limit. The websurfer randomly clicks an random outgoing link with probability $\alpha$ and jumps to a random node with probability $1 - \alpha$.
$$x_i = \alpha \sum_{j=0}^{n-1} \underbrace{\frac{A_{ij}}{k_j} x_j}_{\text{click from $j$}} + \underbrace{\frac{1 - \alpha}{n}}_{\text{random jump}}$$

Considering the probability of going $k$ steps without teleporting
$$x_i = \sum_{k=0}^{\infty} (1 - \alpha) \alpha^k \, \left(\frac{1}{n} \sum_{j=0}^{n-1} (W^k)_{ji}\right) = \frac{1-\alpha}{n} \sum_{j=0}^{n-1} \left((I - \alpha W)^{-1}\right)_{ji}$$
where $W_{ij} = A_{ij}/k_j$ is the walk matrix.

\textbf{Betweenness centrality:} the importance of a node is proportional to the number of shortest paths that pass through it.
$$
b_i = \sum_{j,k} \frac{\overbrace{s_{jki}}^{\text{\# shortest paths through } i}}{\underbrace{s_{jk}}_{\text{\# shortest paths between $j$ and $k$}}}
$$

For a fully connected graph, $b_i = 0$ for all nodes $i$. For two separate fully connected graphs, each with $n$ nodes, connected by a single edge, the nodes on either side of the edge have $b_i = (n-1)n$ and all other nodes have $b_i = 0$.

\subsection{Clustering Coefficient}
A triangle forms when $$A_{ij} A_{jk} A_{ik} = 1$$

\textbf{Local clustering coefficient:} 
$$c_i = \frac{\text{\# triangles with node $i$}}{\text{\# maximum possible triangles with node $i$}} = \frac{\sum_{j,k} A_{ij} A_{jk} A_{ik}}{k_i(k_i - 1)}$$

A clustering coefficient for the whole network can be defined as the average of $c_i$, $c = \frac{1}{n} \sum_i c_i$.

\textbf{Global clustering coefficient (transitivity):}
$$c = \frac{\text{\# triangles}}{\text{\# maximum possible triangles}} = \frac{\sum_{i,j,k} A_{ij} A_{jk} A_{ik}}{\sum_i k_i(k_i - 1)}$$

\subsection{Assortativity and Modularity}
Motivation: `homophily' – nodes with similar properties tend to connect to each other.

\textbf{Assortativity:} a measure of the extent to which nodes' properties are correlated across edges. The assortativity score is the Pearson correlation:
$$
r = \frac{\text{cov}(x_i, x_j)}{\sqrt{\text{var}(x_i)\text{var}(x_j)}} = \frac{\frac{1}{2m} \sum_{i,j} A_{ij} x_i x_j \, - \, \left(\frac{1}{2m}\sum_i k_i x_i\right)^2}{\frac{1}{2m}\sum_i k_i x_i^2 \, - \, \left(\frac{1}{2m}\sum_i k_i x_i\right)^2}
$$

If $x_i = x_j$ whenever $A_{ij} = 1$, then $r = 1$ and corresponds to perfect assortativity. 
Conversely, if $x_i = -x_j$ whenever $A_{ij} = 1$, then $r = -1$ and corresponds to perfect disassortativity.

\textbf{Modularity:} generalises assortativity for categories, beyond numerical quantities. Notations: $g_i$ = category of node $i$, $\delta_{g_i,r} = 1$ if $g_i = r$ and 0 otherwise.
The modularity score quantifies how likely nodes in group $r$ are connect to others in the same group:
$$
Q = \sum_r \left( \frac{1}{2m}\sum_{i,j} A_{ij} \underbrace{\delta_{g_i,r} \delta_{g_j,r}}_{\delta_{g_i, g_j}} - \left(\frac{1}{2m}\sum_i k_i \delta_{g_i,r}\right)^2 \right)
= \frac{1}{2m} \sum_{i,j} \left(A_{ij} - \frac{k_i k_j}{2m} \right) \delta_{g_i, g_j}
$$

\section{Spectral Methods and Interpolation}
\subsection{Linear Interpolation}
Interpolate $y_i$ for node $i$ using
$$
y_i =
\begin{cases}
y_i & \text{if node $i$ is observed} \\
\frac{1}{k_i}\sum_j A_{ij} y_j & \text{otherwise}
\end{cases}
$$

Define $W_{ij} = A_{ij}/k_i,
    y_i = \sum_j W_{ij} y_j
    = \sum_{j:\text{obs.}} W_{ij} y_j + \sum_{j:\text{unobs.}} W_{ij} y_j
    = b_i + \sum_{j:\text{unobs.}} W_{ij} y_j$.

Define $W_{ij}^{(u)} = W_{ij}$ for unobserved nodes $i$ and $j$ gives
\begin{align*}
\mathbf{y} &= \mathbf{b} + W^{(u)} \mathbf{y} \\
\mathbf{y} &= \left(I - W^{(u)}\right)^{-1} \mathbf{b}
\end{align*}

\subsection{Graph Fourier Transform}
Laplacian operator $\mathcal{L}$ is defined as:
$$
\mathcal{L} f = \frac{\mathrm{d}^2 f}{\mathrm{d} x^2} = \lim_{h \to 0} \frac{f(x + h) - 2f(x) + f(x - h)}{h^2}
$$

On a graph, $h = 1$ and $f(x) - f(x + h)$ corresponds to $y_i - y_j$, yielding:
$$
(\mathcal{L} \mathbf{y})_i = \sum_j A_{ij} (y_i - y_j) = k_i y_i - \sum_j A_{ij} y_j
$$
The Laplacian matrix is $\mathcal{L} = D - A$, where $D$ is the diagonal degree matrix with $D_{ii} = k_i$ and $D_{ij} = 0$ for $i \neq j$.
The corresponding eigenvalue equation is $\mathcal{L} \mathbf{v} = \lambda \mathbf{v}$.

The signals can be written as a series in the eigenvectors of $\mathcal{L}$:
$$\mathbf{y} = \sum_{\alpha=0}^{n-1} c_\alpha \mathbf{v}_\alpha \text{ with } c_\alpha = \mathbf{v}_\alpha^\top \mathbf{y}$$

Facts:
\begin{enumerate}
    \item The Laplacian matrix $\mathcal{L}$ is positive semi-definite, i.e. $\lambda_i \geq 0$ for all $i$.
    \item The constant vector $\mathbf{v}_0 = (1, 1, \ldots, 1)^\top$ is an eigenvector of $\mathcal{L}$ with eigenvalue $\lambda_0 = 0$.
    \item The eigenvalue $\lambda_1 = 0$ if and only if the graph is not connected.
\end{enumerate}

\subsection{Spanning Trees}
\textbf{Tree:} a network without any cycles.

\textbf{Spanning tree:} a subset of edges that forms a tree on its nodes.

\textbf{Deletion contraction relation:} the number of spanning trees in a graph $G$ is equal to the number of spanning trees in $G$ with an edge $e$ deleted ($G \backslash e$) plus the number of spanning trees in $G$ with the same edge $e$ contracted ($G / e$).
$$
t(G) = t(G \backslash e) + t(G / e)
$$

Facts:
\begin{enumerate}
    \item The number of spanning trees for a graph with Laplacian $\mathcal{L}$ is equal to $|\mathcal{L}^{(i,i)}|$ (determinant of the matrix obtained by removing the $i$-th row and column of $\mathcal{L}$), for any $i$.
    \item The number of spanning trees of a graph is $\frac{1}{n} \prod_{i=1}^{n-1} \lambda_i$.
\end{enumerate}

\subsection{Normalised Laplacian}
\textbf{Random walk Laplacian:} $\mathcal{L}^{\text{(rw)}} = I - W$ with $W_{ij} = A_{ij}/k_i$.

\textbf{Symmetric normalised Laplacian:}
$$
\mathcal{L}_{ij}^{\text{(sn)}} = 
\begin{cases}
1 & \text{if } i = j \\
-\frac{A_{ij}}{\sqrt{k_i k_j}} &\text{otherwise}
\end{cases}
$$

Relationship between different Laplacians:
\begin{align*}
\mathcal{L} &= D - A \\
\mathcal{L}^{\text{(rw)}} &= D^{-1} \mathcal{L} \\
\mathcal{L}^{\text{(sn)}} &= D^{-\frac{1}{2}} \mathcal{L} D^{-\frac{1}{2}}
\end{align*}

\subsection{Clustering}
Find an $\mathbf{x}$ that minimises
$$\sum_{i,j} A_{ij} (x_i - x_j)^2 = \mathbf{x}^\top \mathcal{L} \mathbf{x}$$
and assign group 1 if $x_i \geq 0$ and group 2 if $x_i < 0$, subject to constraints:
\begin{enumerate}
    \item $\sum_i x_i = 0$ -- to avoid nodes to be in one group.
    \item $\sum_i x_i^2 = 1$ -- to ensure variation.
\end{enumerate}

This yields $\mathbf{x}^\top \mathcal{L} \mathbf{x} = \sum_{j=1}^{n-1} c_j^2 \lambda_j$, which is minimised when $\mathbf{x} = \mathbf{v}_1$, the second eigenvector of $\mathcal{L}$.

\section{Random Graph Models}
\textbf{The random graph $G(n,p)$:} a graph with $n$ nodes where each pair of nodes is connected with probability $p$. The adjacency matrix follows $A_{ij} \sim \text{Bernoulli}(p)$.
For an undirected graph (symmetric adjacency matrix, only consider each edge once)
$$
P(A) = \prod_{i < j} p^{A_{ij}} (1-p)^{1 - A_{ij}}
$$

The probability of a node to be degree $k$ is $p_k = \binom{n-1}{k} p^k (1-p)^{n-1-k}$

Given an adjacency matrix, the maximum likelihood estimate of $p$ is $\hat{p} = \frac{2m}{n(n-1)}$.

\textbf{Poisson random graph:} a random graph with $p = \frac{\lambda}{n}$. With $n \to \infty$, the degree distribution converges to a Poisson distribution $p_k = \frac{\lambda^k e^{-\lambda}}{k!}$.

\textbf{Giant component:} as $n \to \infty$, the size of the largest component would go to infinity -- the giant component.

Facts:
\begin{enumerate}
    \item There is at most one giant component in a random graph.
    \item For $p > 1/n$, there exists a giant component.
\end{enumerate}

\textbf{Triangles:} In a Poisson random graph, the total number of triangles is
$$
\mathbb{E}\left[\sum_{i<j<k} A_{ij} A_{jk} A_{ik}\right] = \binom{n}{3} p^3 \rightarrow \frac{\lambda^3}{6} \text{ as } n \to \infty
$$
The probability for a randomly chosen node to have a triangle is 0.

\textbf{Squares:} In a Poisson random graph, the total number of cycles of length 4 (squares) is
$$
\mathbb{E}[\text{\#4-cycles}] = 3\binom{n}{4} p^4 (1-p)^2 \rightarrow \frac{\lambda^4}{8} \text{ as } n \to \infty
$$

\textbf{General motifs:} For a motif $G'$ with $s$ nodes and $k$ edges, the probability of seeing a $G'$ in a random graph is proportional to $p^k (1-p)^{\binom{s}{2} - k}  = \mathcal{O}(n^{-k})$.
$$
\mathbb{E}[\# G'] = \binom{n}{s} \underbrace{Z(G')}_{\text{a property of $G'$}} p^k (1-p)^{\binom{s}{2} - k} = \mathcal{O}(n^{s-k}) \rightarrow
\begin{cases}
0 & \text{if } k > s \\
\text{const.} & \text{if } k = s
\end{cases}
$$

Implication: as $n \to \infty$, a randomly chosen node is unlikely to be part of any cycles, hence the random graph is locally tree-like.

\textbf{The small components:} any node not in the unique gian component must be in a small component. The small components are distributed as 
$$
\pi{s} = \frac{(s\lambda)^{s-1} e^{-s\lambda}}{s!}
$$

\textbf{Unrealistic properties of the random graph:}
\begin{enumerate}
    \item Locally tree-like and does not contain a high density of short cycles.
    \item Has a Poisson degree distribution, rather than a heavy-tailed degree distribution.
\end{enumerate}

\subsection{Watts-Strogatz Model}
To sample from a Watts-Strogatz model:
\begin{enumerate}
    \item Create a cycle on $n$ nodes.
    \item Connect each node to its $k$ (should be even) nearest neighbors on the cycle. Each node is in $\frac{3k(k-2)}{8}$ traingles and the average shortest path length is $l = \frac{n+k-1}{2k}$.
    \item For every edge $(i,j)$, rewire it to $(i, j')$ for a randomly selected node $j'$ with probability $p$.
\end{enumerate}

Eventually, every node will have at least $k/2$ edges. The number of extra edges is approximately Poisson distributed with mean $k/2$. The large components fo these networks have low diameter.

Fact: For any fixed $p > 0$, the Watts-Strogatz model generates the small-world property, i.e. the distance between nodes is $\mathcal{O}(\log n)$.

\subsection{Preferential Attachment Model (Barabási-Albert Model)}
In a preferential attachment model, nodes are added to the network one at a time. As a new node arrives, it preferentially attaches to $c$ nodes with a large number of existing edges, i.e. proportional to the degree.

Given $n$ nodes in the network, the $(n+1)$th node selects node $i$ to connect to with probability $$\frac{k_i}{\sum_j k_j} = \frac{k_i}{2cn}$$

Fact: The preferential model has a power-law degree distribution $p_k \sim k^{-3}$.

Note: The distribution is heavy-tailed. Mean degree $\sum_k k p_k = 2c$, the variance of degree is $\sum_k (k - 2c)^2 p_k = \infty$ as $n \to \infty$.

\subsection{Graphs with Arbitrary Degree Distribution}
\textbf{Configuration model:} For each node $i$, assign it a degree $k_i$ drawn from the degree distribution $p_k$. Create $k_i$ stubs (half-edges) and connect stubs at random to form edges. The resultant graph is a multigraph with exactly the degree sequence $\{k_i\}$.

\textbf{Excess degree distribution $q_k$:} probability that a node reached by following a random edge has $k$ other edges. This is proportional to the number of edges that can be followed to reach the node, i.e. $q_k \propto (k + 1)p_{k+1}$, normalised by $1/\mathbb{E}[k]$.

The generating function for $q_k$ is 
$$
g_q(z) = \frac{g'_p(z)}{g'_p(1)} = \frac{\sum_k k p_k z^{k-1}}{\mathbb{E}[k]} = \sum_{k=0}^\infty q_k z^k
$$

Define $u$ = probability a random not is not in the giant component, and $w$ = probability a random edge does not lead to the giant component, then $u$ and $w$ satisfy:
\begin{align*}
    u &= \sum_k p_k w^k = g_p(w) \\
    w &= \sum_k q_k w^k = g_q(w) = \frac{g'_p(w)}{g'_p(1)}
\end{align*}

Fact: A random network with degree distribution $p_k$ will have a giant component if $\mathbb{E}[k^2] > 2\mathbb{E}[k]$.

\subsection{Graphs with Community Structure}
\textbf{Stochastic block model:} a model that accounts for varying densities within and between groups. Each node $i$ is placed in one of $K$ groups at random, i.e. $g_i \in \{1, \dots, K\}$. 
The adjacency matrix is generatated from $A_{ij} \sim \text{Bernoulli}(\omega_{g_i g_j})$, where $\omega_{g_i g_j}$ is the probability of an edge between nodes in groups $g_i$ and $g_j$. 
With group assignment $\boldsymbol{g}$, the probability for an adjacency matrix is
$$
P(A| \boldsymbol{g}, \omega) = \prod_{i < j} \omega_{g_i g_j}^{A_{ij}} (1 - \omega_{g_i g_j})^{1 - A_{ij}}
$$

\section{Community Detection}
Goal: find statistically related nodes in a network.

\subsection{Clustering}
To recover the communities $\boldsymbol{g}$ given a network $A$, apply Bayes' law:
$$
P(\boldsymbol{g} | A, \omega) = \frac{P(A | \boldsymbol{g}, \omega) P(\boldsymbol{g})}{P(A | \omega)} = \frac{1}{Z}\prod_{i < j} \omega_{g_i g_j}^{A_{ij}} (1 - \omega_{g_i g_j})^{1 - A_{ij}} P(\boldsymbol{g})
$$

Compute the posterior probability $Q_{ir}$ that node $i$ is in group $r$ using naive Bayes:
\begin{align*}
    Q_{ir} &= P(g_i = r | A, \omega) = \sum_{\boldsymbol{g}\backslash_i} P(g_i = r, \boldsymbol{g}\backslash_i | A, \omega) \\
    &\propto \sum_{\boldsymbol{g}\backslash_i} \left(\prod_j \omega_{r g_j}^{A_{ij}} (1 - \omega_{r g_j})^{1 - A_{ij}}\right) \underbrace{\left(\prod_{j < k, \neq i} \omega_{g_j g_k}^{A_{jk}} (1 - \omega_{g_j g_k})^{1 - A_{jk}}\right)}_{\text{posterior when $i$ is removed}}
\end{align*}

Assuming (1) removing a single node $i$ would not change the posterior for the rest and (2) the posterior term factorises (true in appropriate limits) gives:
\begin{align*}
    Q_{ir} &\propto \sum_{\boldsymbol{g}\backslash_i} \prod_j Q_{jg_j} \omega_{r g_j}^{A_{ij}} (1 - \omega_{r g_j})^{1 - A_{ij}} \\
    &= \prod_j \sum_s \omega_{r s}^{A_{ij}} (1 - \omega_{r s})^{1 - A_{ij}} Q_{js}
\end{align*}

Given $Q_{ir}$, the minimum error is achieved by setting $\hat{g}_i = \arg\max_r Q_{ir}$.

\subsection{Parameter Estimation}
Find $\omega$ that maximises the likelihood of the observed network $A$:
\begin{align*}
\log P(A | \omega) &= \log \sum_{\boldsymbol{g}} P(A , \boldsymbol{g} | \omega) \\
&\geq \sum_{\boldsymbol{g}} Q(\boldsymbol{g}) \log \frac{P(A, \boldsymbol{g} | \omega)}{Q(\boldsymbol{g})} \quad\text{(Jensen's inequality)} \\
\end{align*}
Setting $Q(\boldsymbol{g}) = P(\boldsymbol{g} | A, \omega)$ equates the LHS and the RHS, and the optimal $\omega$ is
$$
w_{rs} = \frac{\sum_{i,j} A_{ij} Q_{ir} Q_{js}}{\sum_{i,j} Q_{ir} Q_{js}}
$$

Use EM algorithm to iteratively update $Q_{ir}$ and $\omega$ until convergence:
\begin{enumerate}
    \item E-step: update $Q_{ir}$ using the current $\omega$.
    \item M-step: update $\omega$ using the current $Q_{ir}$.
\end{enumerate}

\section*{Notations}
\renewcommand{\arraystretch}{1.2}
\begin{center}
\begin{tabular}{|c|l|}
\hline
\textbf{Symbol} & \textbf{Definition} \\ \hline
$G$ & A network – $G$ for graph. \\ \hline
$n$ & Number of nodes in a network. \\ \hline
$m$ & Number of edges in a network. \\ \hline
$k$ & Mean degree, equal to $2m/n$. \\ \hline
$k_i$ & Degree of node $i$. \\ \hline
$A$ & Adjacency matrix. \\ \hline
$\mathcal{L}$ & Laplacian matrix. \\ \hline
$p_k$ & Degree distribution. \\ \hline
$q_k$ & Excess degree distribution, $q_k \propto (k + 1)p_{k+1}$. \\ \hline
$G(n,p)$ & The random graph. \\ \hline
$N_i$ & The set of neighbors of node $i$. \\ \hline
\end{tabular}
\end{center}

\end{document}